% Mục 6.4 (bản rút gọn). Rút từ sections/results_cross.tex; nguồn từng số ghi ở đó (results/<máy>/*.csv,
% machine.txt; scripts/report_numbers/rcross_*.py). Chỉ giữ các cài đặt có trên Hình 7--9 (CROSS_MAIN của plot.py).
\subsubsection{RQ6: Kết luận trên máy khác}

\begin{rqanswer}\textbf{Trả lời RQ6.} Một phần. Thứ hạng chính giữ trên cả ba máy. TTAS nhanh hơn TAS,
fetch-and-add nhanh hơn vòng CAS từ 4 thread, \PCname{sr_bo_yield} tốt nhất khi quá tải, và backoff bằng
ngủ tuần tự hoá các thread. Giữa hai máy Apple, điểm gãy dời theo số P-core và điểm sụp dời theo số core.
Trên \PCname{LAP1184} (x86-64, Linux trong máy ảo WSL2), \texttt{seq\_cst} không còn tốn thêm gì, barrier
spin chậm 17 lần ngay ở 12 thread (bằng số luồng phần cứng), \texttt{sleep\_for} với yêu cầu 100~ns ngủ trung
vị khoảng 62~\textmu s, và MCS quá tải cho hai cực thời gian.\end{rqanswer}

Hai máy Apple khác nhau rõ nhất ở số P-core và số core, nên chúng cho phép so tác động của hai con số đó.
\PCname{LAP1184} khác \PCname{apple-m5} ở kiến trúc, memory model, compiler, hệ điều hành, và chạy trong
máy ảo. Hình phần lock (Hình~\ref{fig:cross:lock}) chia thời gian cho cận dưới của chính máy đó, nên khử
được khác biệt tốc độ core. Hình barrier và counter (Hình~\ref{fig:cross:barrier}
và~\ref{fig:cross:counter}) vẽ số đo thô, nên chỉ so được hình dạng đường cong. Trên ô \PCname{LAP1184},
vạch chấm ở 6 là số core máy ảo báo, còn vạch liền ở 12 là số luồng phần cứng.

\PCfigureopt{cross}{lock}

\emph{Bằng chứng.} Lock (Hình~\ref{fig:cross:lock}). Trên \PCname{C4M7Y670NT}, TAS chậm 1,15 lần cận dưới
ở 6 thread và 1,53 ở 8 thread, MCS 1,23 rồi 2,58, tức chỗ gãy nằm ở 6 lên 8 thay vì 4 lên 6. MCS chậm
4,99 lần ở 12 thread (bằng số core) và 2\,070--5\,920 lần khi vượt. Trên \PCname{LAP1184}, tới 12 thread
TAS, TTAS và \PCname{omp_lock} chạy trong 1,36 lần cận dưới, không gãy ở 6 hay 10, và ở 24 thread ở mức
1,7--2,2. MCS chậm 1,85 lần ở 10 thread và 3,13 lần ở 12 thread. TAS chậm hơn TTAS 1,30 lần ở 24 thread.
\PCname{std_mutex} chậm 2,0 lần cận dưới ngay ở 2 thread.
% nguồn: C4M7Y670NT locktime N=1009 OW=0: tas 1,15/1,53; mcs 1,23/2,58/4,99, 2068..5920. LAP1184 locktime: 12t tas
% 1,36, ttas 1,26, omp_lock 1,27; 24t tas 2,24, ttas 1,72, omp_lock 1,79; tas/ttas 24t 1,30; std_mutex 2t 1,96;
% mcs 10t 1,85, 12t 3,13.

\PCfigureopt{cross}{barrier}

Barrier (Hình~\ref{fig:cross:barrier}). Trên \PCname{C4M7Y670NT}, mỗi thread thêm vào làm một episode của
\PCname{sr} dài thêm 36--39~ns tới 6 thread, rồi 139~ns từ 6 lên 8 thread. \PCname{sr_block} và
\PCname{std_barrier} tụt 6,6 và 5,9 lần ở cùng chỗ. \PCname{sr} vẫn đạt 1,35 triệu episode mỗi giây ở 12
thread và chỉ sụp khi vượt số core (184 ở 16 thread). Trên \PCname{LAP1184}, \PCname{sr} ở 1 thread chỉ
đạt 128 triệu (\PCname{apple-m5} 589 triệu). Từ 4 tới 10 thread, \PCname{sr} dài thêm 77--101~ns mỗi
thread, không gãy ở 6 hay 10, nhưng ở 12 thread nó chậm 17 lần so với 10 thread. \PCname{omp_barrier}
chậm hơn \PCname{sr} từ 1 tới 10 thread, vượt \PCname{sr} ở 12 thread (1,03 triệu), rồi sụp còn 4,25
nghìn ở 16 thread. Các barrier chờ chặn rơi xuống 18--26 nghìn ngay ở 3 thread. Khi quá tải,
\PCname{sr_bo_yield} nhanh nhất trên cả ba máy.
% nguồn: C4M7Y670NT barrier sr: (137-66)/2=36, (215-137)/2=39, (493-215)/2=139; sr_block 6,6, std_barrier 5,9;
% sr 12t 1,35M, 16t 184. LAP1184 barrier sr 1t 128M (apple-m5 589M); 77..101 ns mỗi thread; 1,35M/78,2k=17,3;
% omp_barrier 12t 1,03M, 16t 4,25k; chờ chặn 3t 18,2k-25,9k.

\PCfigureopt{cross}{counter}

Counter (Hình~\ref{fig:cross:counter}). Trên \PCname{C4M7Y670NT}, một lần tăng của \PCname{faa} tốn
5,8--8,1~ns ở 4--6 thread và 21--38~ns từ 8 thread. \PCname{cas} nhanh hơn \PCname{faa} ở 2 thread (0,76
lần, \PCname{apple-m5} 0,72), rồi chậm hơn 3,6--11 lần từ 4 thread. Counter dùng mutex nhanh hơn
\PCname{faa} từ 8 thread (\PCname{lock_std_mutex} 17,8 so với 27,0~ns), nhưng ở 3--4 thread ba lock thư
viện chậm hơn \PCname{faa} 6,6 tới 25 lần. Trên \PCname{LAP1184}, \PCname{cas} chậm hơn \PCname{faa} ngay
từ 2 thread (1,49 lần) và 2,9 tới 4,3 lần từ 4 thread, dù số lần hỏng ở 8 thread gần bằng
\PCname{apple-m5} (1,62 và 1,59). \PCname{lock_mcs} tăng từ 159 lên 918~ns mỗi lần tăng khi đi từ 8 lên
10 thread, ở cả bảy lần chạy. Từ 12 thread, \PCname{lock_omp_lock} chậm hơn \PCname{lock_tas} 1,2--1,4
lần.
% nguồn: C4M7Y670NT counter OW=0 ns/op (rcross_m3.py): faa 4t 5,8, 6t 8,1, 8..24t 21,6-38,2; cas/faa 2t 0,76,
% 4..24t 3,56-11,1; std_mutex 8t 17,8, faa 27,0; mutex/faa 3..4t 6,6..25. LAP1184: cas/faa 2t 1,49, 4..24t
% 2,87..4,29; casfail 8t 1,62; lock_mcs 8t 159,1, 10t 917,8; lock_omp_lock/lock_tas 12..24t 1,20..1,39.

\emph{Vì sao: hai máy Apple.} Chỗ gãy của spinlock, của độ dốc barrier spin và của các barrier
quay-rồi-ngủ nằm ngay sau số P-core trên cả hai máy. Điều này phù hợp với giải thích ở RQ2 và RQ4, theo đó
thread vượt số P-core chạy trên E-core chậm hơn, và cache line phải đi giữa hai cụm core. \PCname{faa}
trên \PCname{C4M7Y670NT} rất rẻ trong cụm 6 P-core dùng chung L2, và tốn gấp 2,7 tới 4,7 lần khi có
thread ở cụm E-core. Ở 4--6 thread, \PCname{cas} tốn gấp 7,4--11 lần \PCname{faa}, lớn hơn khoảng 3 lần
mức mô hình $1+f$ của RQ5. Khoảng chênh này phù hợp với giả thuyết của RQ5 rằng \texttt{ldadd} được thực
hiện ở cache dùng chung, còn CAS phải kéo line về L1. MCS và barrier spin sụp khi vượt số core, vì thread
chờ phải đợi một lượt lập lịch. Các chi phí do hệ điều hành quyết định thì bằng nhau, ví dụ
\texttt{sleep\_for} với yêu cầu 100~ns ngủ trung vị 2,67~\textmu s trên cả hai máy.
% nguồn: faa C4M7Y670NT 2,7..4,7; cas/faa 4t 7,4, 6t 11,1; 7,4/2,42=3,1, 11,1/3,82=2,9 (rcross_m3.py).

\emph{Vì sao: \PCname{LAP1184}.} Thứ nhất, một lời gọi \texttt{is\_prime(1009)} tốn 1\,640~ns, gấp 3,2
lần \PCname{apple-m5}. Lấy thời gian mỗi critical section ở 20 thread trừ đi $t$ được chi phí thêm cho
mỗi lần qua lock. Với TAS, chi phí này gần bằng nhau trên hai máy (1,34 và 1,27~\textmu s), nên chia cho
critical section dài hơn nó cho độ chậm nhỏ hơn. Với mutex, chi phí thêm gấp 2,2--2,3 lần
\PCname{apple-m5}, phù hợp với việc đánh thức qua futex (lời gọi hệ thống của Linux để thread ngủ và được
đánh thức trên một biến) trong máy ảo đắt hơn. Nhóm cho rằng Windows đặt
12 CPU ảo lên P-core và E-core mà Linux không biết, nên đường spinlock không gãy ở một số thread cố định.
% nguồn: lockunit t N=1009 OW=0: 1\,640 / 511 ns; TAS 20t 1,34 / 1,27 µs; std_mutex 3,37..3,51 / 1,50 = 2,25..2,34.

Thứ hai, trên x86-64 mọi phép đọc-sửa-ghi atomic là lệnh có tiền tố \texttt{lock}, vốn đã là một rào
chắn đầy đủ, nên \texttt{seq\_cst} không đổi lệnh máy. \PCname{cas_seqcst} bằng 0,75 tới 1,14 lần
\PCname{cas} và \PCname{faa_seqcst} bằng 0,90 tới 1,13 lần \PCname{faa}. \texttt{lock xadd} và \texttt{lock cmpxchg} đều kéo line về L1 của core gọi, nên không có
hiệu ứng \texttt{ldadd} ở xa, phù hợp với việc \PCname{cas} không thắng \PCname{faa} ở 2 thread. Lệnh
\texttt{pause} tốn trung vị 51~ns, gấp 4 lần \texttt{isb}, nên mỗi lần CAS hỏng chờ lâu hơn.
% nguồn: LAP1184 countertime OW=0 cas_seqcst/cas 0,75..1,14; sleep_latency.csv: pause 50,7 ns, isb 12,1 ns.

Thứ ba, với yêu cầu 100~ns, \texttt{sleep\_for} ngủ trung vị 61,7~\textmu s, gấp 23 lần hai máy Apple.
Nhóm cho rằng phần lớn khoảng này là timer slack mặc định 50~\textmu s của Linux (khoảng trễ kernel được
phép thêm khi đánh thức để gộp các lần hẹn giờ) cộng độ trễ bộ hẹn giờ của máy ảo. Vì vậy, ở
$\mathrm{OW} = 3$, \PCname{tas_bo_sleep} chậm hơn TAS 1,5--2,0 lần ở 3--16 thread, trong khi trên
\PCname{apple-m5} nó nhanh hơn TAS ở 20 thread. Thứ tư, barrier chờ chặn tốn 39--74~\textmu s mỗi episode
ở 3--4 thread. Nhóm cho rằng \texttt{atomic::wait} của libstdc++ chỉ quay rất ngắn trước khi gọi futex.
Đây là giả thuyết.
% nguồn: LAP1184 sleep_latency.csv sleep_for 100 ns trung vị 61\,700 ns; 61\,700/2\,670=23; locktime-ow3
% tas_bo_sleep/tas 3..16t 1,54-2,04; chờ chặn 3..4t 39..74 µs.

Thứ năm, \PCname{omp_barrier} cũng quay ở 12 thread mà vẫn đạt khoảng 1~\textmu s mỗi episode, nên riêng
việc 12 thread cùng quay không làm barrier chậm. Nhóm cho rằng điểm khác là ai nhường CPU. Khi một thread
của \PCname{sr} mất CPU, 11 thread còn lại vẫn quay và giữ CPU, nên cả barrier chờ tới khi thread đó chạy
lại. libgomp quay một thời gian rồi ngủ bằng futex (theo tài liệu \texttt{GOMP\_SPINCOUNT}), nên trả CPU
lại. \PCname{sr_bo_yield} cũng nhường CPU và đạt 111 nghìn episode mỗi giây ở 12 thread, hơn \PCname{sr}
(78,2 nghìn). MCS chậm 3,13 lần cận dưới ở 12 thread, phù hợp với cùng cơ chế.
% nguồn: LAP1184 barrier 12t: omp_barrier 1,03M (0,97 µs), sr 78,2k, sr_bo_yield 111k.

\emph{Độ vững.} Khi quá tải, MCS cho hai cực thời gian. Trên hai máy Apple, gần như mọi lần chạy quá tải
đều chậm (41 trên 42 lần). Trên \PCname{LAP1184}, trong 21 lần chạy ở 16, 20 và 24 thread, 14 lần xong
trong 2,7--7,1~ms và 7 lần mất 1,06--4,92~s. Trung vị rơi vào cực nhanh ở 16 và 24 thread, cực chậm ở 20
thread, rồi được nhân lên 1\,600--2\,400 lần khi quy đổi. Vì vậy điểm thấp ở 24 thread trên
Hình~\ref{fig:cross:lock} không có nghĩa MCS hồi phục. Từ 12 thread, \PCname{cas} và \PCname{cas_seqcst}
trên \PCname{C4M7Y670NT} cùng chia hai cực, nên tỉ số trung vị của chúng (0,69 tới 1,68) chỉ phản ánh
trung vị rơi vào cực nào. Từ mức thread ngay dưới lên số core, \PCname{sr} trên \PCname{apple-m5} chậm
đi 2,7 lần (8 lên 10 thread), còn trên \PCname{C4M7Y670NT} chỉ 1,17 lần (10 lên 12 thread). Điều này phù
hợp với tải nền cao hơn của \PCname{apple-m5} (trung vị load average 1 phút trong phần barrier là 26 so
với 18).
% nguồn: rcross_mcs.py; LAP1184 lock.csv mcs N=1009 OW=0 16/20/24t; C4M7Y670NT counter.csv cas, cas_seqcst
% 12..24t; sr C4M7Y670NT 1,17, apple-m5 2,68; load1 trung vị barrier.csv 26,5 và 18,3.
